\chapter{Dérivation et convexité}\label{ch:g12:deriv}

La \omterm{def:g12:deriv:derivative}{dérivée} mesure le taux de variation instantané d'une \omterm{def:g11:func:function}{fonction}~; c'est le
\omterm{def:g10:reffunc:affine}{coefficient directeur} de la \omterm{def:g12:deriv:tangent}{tangente} à son \omterm{def:g10:functions:graph}{graphique}. Ce chapitre revoit et
étend les règles de dérivation, relie le signe de $f'$ aux variations de $f$,
et introduit la \omterm{def:g12:deriv:convex}{convexité}, gouvernée par la \omterm{def:g12:deriv:derivative}{dérivée} seconde.

\section{La dérivée}

\begin{definition}[Dérivée en un point]\label{def:g12:deriv:derivative}
Soit $f$ définie sur un \omterm{def:g10:numbers:interval}{intervalle} $I$ et $a \in I$. La \omterm{def:g11:func:function}{fonction} $f$ est
\emph{dérivable en $a$}\index{dérivée}\index{dérivable} si le taux d'accroissement
\[
\frac{f(a+h) - f(a)}{h}
\]
a une limite finie lorsque $h \to 0$. Cette limite s'appelle la
\emph{dérivée de $f$ en $a$} et se note $f'(a)$. Si $f$ est dérivable en
tout point de $I$, la \omterm{def:g11:func:function}{fonction} $f' \colon x \mapsto f'(x)$ est la
\emph{dérivée} de $f$.
\end{definition}

\begin{definition}[Tangente]\label{def:g12:deriv:tangent}
Si $f$ est \omterm{def:g12:deriv:derivative}{dérivable} en $a$, la \emph{tangente}\index{tangente} à la courbe
de $f$ au point $(a, f(a))$ est la droite d'\omterm{def:g10:algebra:equation}{équation}
\[
y = f(a) + f'(a)(x - a).
\]
\end{definition}

\begin{omfigure}
\begin{tikzpicture}
\begin{axis}[omaxis, clip=false,
  width=9.6cm, height=6.8cm,
  xmin=-0.4, xmax=3.4, ymin=-0.4, ymax=4.3,
  xtick={1,2.5}, xticklabels={$a$,$a+h$}, ytick=\empty]
  \addplot[omProp!40, thick, domain=0.2:2.7] {1 + 1.75*(x-1)};
  \addplot[omProp!70, thick, domain=0.07:2.95] {1 + 1.5*(x-1)};
  \addplot[omProp, thick, domain=-0.12:3.3] {1 + 1.25*(x-1)};
  \addplot[omThm, thick, domain=-0.4:3.3] {x};
  \addplot[omDef, thick, domain=-0.4:2.72] {x^2/2 + 0.5};
  \draw[black, dashed, thin] (axis cs:1,1) -- (axis cs:2.5,1)
    node[midway, below, font=\small] {$h$};
  \draw[black, dashed, thin] (axis cs:2.5,1) -- (axis cs:2.5,3.625)
    node[midway, right, font=\small] {$f(a+h)-f(a)$};
  \fill (axis cs:2.5,3.625) circle (1.5pt);
  \fill (axis cs:2,2.5) circle (1.5pt);
  \fill (axis cs:1.5,1.625) circle (1.5pt);
  \fill (axis cs:1,1) circle (1.5pt)
    node[above left, font=\small] {$A$};
  \node[omThm, font=\small, anchor=west] at (axis cs:3.38,3.3) {tangente};
\end{axis}
\end{tikzpicture}

{\small Lorsque $h$ diminue, la \omterm{def:g11:deriv:rate}{sécante} par $A = (a, f(a))$ et
$(a+h, f(a+h))$ (orange) tourne vers la \omterm{def:g12:deriv:tangent}{tangente} en $A$ (rouge)~: son
\omterm{def:g10:reffunc:affine}{coefficient directeur} $\frac{f(a+h)-f(a)}{h}$ tend vers $f'(a)$.}
\end{omfigure}

\begin{proposition}\label{prop:g12:deriv:diffcont}
Une \omterm{def:g11:func:function}{fonction} \omterm{def:g12:deriv:derivative}{dérivable} en $a$ est \omterm{def:g12:limcont:continuity}{continue} en $a$. La réciproque est fausse.
\end{proposition}

\begin{proof}
Pour $h \neq 0$ petit, $f(a+h) - f(a) = h \cdot \frac{f(a+h)-f(a)}{h}$.
Lorsque $h \to 0$, le second membre tend vers $0 \times f'(a) = 0$, donc
$f(a+h) \to f(a)$. Pour la réciproque, $x \mapsto \abs{x}$ est \omterm{def:g12:limcont:continuity}{continue} en
$0$, mais son taux d'accroissement en $0$ égale $\frac{\abs h}{h} = \pm 1$
selon le signe de $h$, et n'a pas de limite.
\end{proof}

\subsection{Règles de dérivation}

\begin{proposition}[Opérations]\label{prop:g12:deriv:operations}
Soient $u, v$ \omterm{def:g12:deriv:derivative}{dérivables} sur $I$ et $\lambda \in \R$. Alors $u + v$,
$\lambda u$, $uv$ sont \omterm{def:g12:deriv:derivative}{dérivables} sur $I$, de même que $\frac{u}{v}$ où
$v \neq 0$, et
\[
(u+v)' = u' + v', \quad (\lambda u)' = \lambda u', \quad
(uv)' = u'v + uv', \quad
\left(\frac{u}{v}\right)' = \frac{u'v - uv'}{v^2}.
\]
\end{proposition}

\begin{proof}[Démonstration de la règle du produit]
Écrire le taux d'accroissement de $uv$ en $a$ comme
\[
\frac{u(a+h)v(a+h) - u(a)v(a)}{h}
= \frac{u(a+h) - u(a)}{h}\,v(a+h) + u(a)\,\frac{v(a+h) - v(a)}{h}.
\]
Lorsque $h \to 0$, $v(a+h) \to v(a)$ par \omterm{def:g12:limcont:continuity}{continuité}
(\cref{prop:g12:deriv:diffcont}), donc le second membre tend vers
$u'(a)v(a) + u(a)v'(a)$. Les autres règles se démontrent de même.
\end{proof}

\begin{theorem}[Règle de la chaîne]\label{thm:g12:deriv:chain}
\index{règle de la chaîne}
Soit $u$ \omterm{def:g12:deriv:derivative}{dérivable} sur $I$ à valeurs dans $J$, et $g$ \omterm{def:g12:deriv:derivative}{dérivable} sur $J$.
Alors $g \circ u$ est \omterm{def:g12:deriv:derivative}{dérivable} sur $I$ et
\[
(g \circ u)' = (g' \circ u) \cdot u' .
\]
En particulier, pour $u$ \omterm{def:g12:deriv:derivative}{dérivable}~:
\[
(u^n)' = n\,u'\,u^{n-1} \ (n \in \Z), \qquad
\left(\sqrt{u}\right)' = \frac{u'}{2\sqrt{u}} \ (u > 0), \qquad
\left(\eu^{u}\right)' = u'\,\eu^{u}.
\]
\end{theorem}

\begin{proof}[Esquisse de démonstration]
Lorsque $u(a+h) \neq u(a)$ pour $h$ petit, écrire
\[
\frac{g(u(a+h)) - g(u(a))}{h}
= \frac{g(u(a+h)) - g(u(a))}{u(a+h) - u(a)} \cdot \frac{u(a+h) - u(a)}{h}.
\]
Lorsque $h \to 0$, $u(a+h) \to u(a)$, donc le premier facteur tend vers
$g'(u(a))$ et le second vers $u'(a)$. (Une démonstration complète doit
traiter le cas où $u(a+h) = u(a)$ pour $h$ arbitrairement petit~; ce point
technique est traité à l'université.)
\end{proof}

\begin{example}
Les \omterm{def:g12:deriv:derivative}{dérivées} usuelles, valables sur les domaines naturels~:
\[
(x^n)' = nx^{n-1}, \quad
\left(\frac{1}{x}\right)' = -\frac{1}{x^2}, \quad
(\sqrt{x})' = \frac{1}{2\sqrt{x}}, \quad
(\cos x)' = -\sin x, \quad
(\sin x)' = \cos x .
\]
La première se prouve par récurrence à partir de la règle du produit, et les
deux dernières dans le \cref{ch:g12:trigo}.
\end{example}

\section{Variations et extrema}

\begin{theorem}[Signe de la dérivée et variations]\label{thm:g12:deriv:variations}
Soit $f$ \omterm{def:g12:deriv:derivative}{dérivable} sur un \omterm{def:g10:numbers:interval}{intervalle} $I$.
\begin{enumerate}
  \item $f$ est \omterm{def:g12:seq:monotonic}{croissante} sur $I$ si et seulement si $f' \geq 0$ sur $I$.
  \item $f$ est \omterm{def:g10:functions:variations}{décroissante} sur $I$ si et seulement si $f' \leq 0$ sur $I$.
  \item Si $f' > 0$ sur $I$ sauf en un nombre fini de points où elle
        s'annule, alors $f$ est strictement \omterm{def:g12:seq:monotonic}{croissante} sur $I$.
\end{enumerate}
\end{theorem}

\begin{proof}[Démonstration partielle]
Si $f$ est \omterm{def:g12:seq:monotonic}{croissante}, tout taux d'accroissement
$\frac{f(a+h)-f(a)}{h}$ est $\geq 0$, et en passant à la limite on obtient
$f'(a) \geq 0$. Les implications réciproques reposent sur le théorème des
accroissements finis, prouvé à l'université~; elles sont \emph{admises à ce
niveau}.
\end{proof}

\begin{definition}[Extremum local]\label{def:g12:deriv:extremum}
$f$ a un \emph{maximum local}\index{extremum} en $a$ si $f(x) \leq f(a)$ pour
tout $x$ voisin de $a$~; les minima locaux se définissent de façon
symétrique.
\end{definition}

\begin{proposition}[Condition du premier ordre]\label{prop:g12:deriv:critical}
Si $f$ est \omterm{def:g12:deriv:derivative}{dérivable} sur un \omterm{def:g10:numbers:interval}{intervalle} \emph{ouvert} $I$ et a un \omterm{def:g12:deriv:extremum}{extremum}
local en $a \in I$, alors $f'(a) = 0$. La réciproque est fausse
($f(x) = x^3$ en $a = 0$).
\end{proposition}

\begin{proof}
Disons que $a$ est un \omterm{def:g12:deriv:extremum}{maximum local}. Pour $h > 0$ petit,
$\frac{f(a+h)-f(a)}{h} \leq 0$, donc en faisant $h \to 0^+$ on obtient
$f'(a) \leq 0$~; pour $h < 0$ le quotient est $\geq 0$, d'où
$f'(a) \geq 0$. Donc $f'(a) = 0$.
\end{proof}

\begin{omfigure}
\begin{tikzpicture}
\begin{axis}[omaxis,
  width=8.4cm, height=5.6cm,
  xmin=-2.5, xmax=2.5, ymin=-4.2, ymax=4.2,
  xtick={-1,1}, ytick={-2,2}]
  \addplot[omProp, thick, domain=-1.8:-0.2] {2};
  \addplot[omProp, thick, domain=0.2:1.8] {-2};
  \addplot[omDef, thick, domain=-2.15:2.15] {x^3 - 3*x};
  \fill (axis cs:-1,2) circle (1.5pt);
  \fill (axis cs:1,-2) circle (1.5pt);
\end{axis}
\end{tikzpicture}

{\small En un \omterm{def:g12:deriv:extremum}{extremum} local à l'intérieur d'un \omterm{def:g10:numbers:interval}{intervalle} ouvert, la
\omterm{def:g12:deriv:tangent}{tangente} (orange) est horizontale~: $f'(a) = 0$. Ici $f(x) = x^3 - 3x$,
avec un \omterm{def:g12:deriv:extremum}{maximum local} en $-1$ et un \omterm{def:g10:functions:extrema}{minimum} local en $1$.}
\end{omfigure}

\begin{method}[Tableau de variations]\label{met:g12:deriv:table}
Pour étudier une \omterm{def:g11:func:function}{fonction} $f$~: déterminer son \omterm{def:g11:func:function}{ensemble de définition} et les
limites au bord~; calculer $f'$ et la \omterm{def:g10:algebra:expand}{factoriser}~; déterminer le signe de
$f'$ sur chaque \omterm{def:g10:numbers:interval}{sous-intervalle}~; tout enregistrer dans un \omterm{def:g10:functions:table}{tableau de
variations}, en marquant les extrema~; en déduire le nombre de solutions de
$f(x) = k$ à l'aide du théorème de la bijection
(\cref{thm:g12:limcont:bijection}).
\end{method}

\section{Convexité}

\begin{definition}[Fonction convexe]\label{def:g12:deriv:convex}
Une \omterm{def:g11:func:function}{fonction} $f$ définie sur un \omterm{def:g10:numbers:interval}{intervalle} $I$ est \emph{convexe}\index{convexité}
sur $I$ si toute corde de son \omterm{def:g10:functions:graph}{graphique} est au-dessus du \omterm{def:g10:functions:graph}{graphique}~: pour
tous $a, b \in I$ et $t \in \intcc{0}{1}$,
\[
f\bigl((1-t)a + t b\bigr) \leq (1-t) f(a) + t f(b).
\]
Elle est \emph{concave}\index{concave} si l'inégalité inverse vaut ($-f$ convexe).
\end{definition}

\begin{omfigure}
\begin{tikzpicture}
\begin{axis}[omaxis,
  width=8.8cm, height=6cm,
  xmin=-1.9, xmax=2.5, ymin=-0.8, ymax=4.3,
  xtick={-1,1.5}, xticklabels={$a$,$b$}, ytick=\empty]
  \addplot[omThm, thick, domain=-1.25:1.75] {0.5*x + 1.5};
  \addplot[omProp, thick, domain=-0.45:2.3] {x - 0.25};
  \addplot[omDef, thick, domain=-1.6:2.0] {x^2};
  \draw[black, dashed, thin] (axis cs:0.25,0.0625) -- (axis cs:0.25,1.625);
  \fill (axis cs:-1,1) circle (1.5pt);
  \fill (axis cs:1.5,2.25) circle (1.5pt);
  \fill (axis cs:0.5,0.25) circle (1.5pt);
  \node[omThm, font=\small, above left] at (axis cs:1.7,2.4) {corde};
  \node[omProp, font=\small, below right] at (axis cs:1.75,1.5) {tangente};
\end{axis}
\end{tikzpicture}

{\small Une \omterm{def:g11:func:function}{fonction} \omterm{def:g12:deriv:convex}{convexe}~: toute corde (rouge) est au-dessus du
\omterm{def:g10:functions:graph}{graphique}, et le \omterm{def:g10:functions:graph}{graphique} est au-dessus de chacune de ses \omterm{def:g12:deriv:tangent}{tangentes}
(orange). Le segment en pointillés montre l'écart entre
$f\bigl((1-t)a+tb\bigr)$ et $(1-t)f(a)+tf(b)$.}
\end{omfigure}

\begin{theorem}[Convexité et dérivées]\label{thm:g12:deriv:convexity}
Soit $f$ \omterm{def:g12:deriv:derivative}{dérivable} sur un \omterm{def:g10:numbers:interval}{intervalle} $I$. Les assertions suivantes sont
équivalentes~:
\begin{enumerate}
  \item $f$ est \omterm{def:g12:deriv:convex}{convexe} sur $I$~;
  \item $f'$ est \omterm{def:g12:seq:monotonic}{croissante} sur $I$~;
  \item le \omterm{def:g10:functions:graph}{graphique} de $f$ est au-dessus de chacune de ses \omterm{def:g12:deriv:tangent}{tangentes}.
\end{enumerate}
Si $f$ est deux fois \omterm{def:g12:deriv:derivative}{dérivable}, elles sont aussi équivalentes à
$f'' \geq 0$ sur $I$.
\end{theorem}

\begin{proof}[Démonstration de (2) $\Rightarrow$ (3)]
Fixer $a \in I$ et poser
$g(x) = f(x) - f(a) - f'(a)(x-a)$, l'écart entre la courbe et la \omterm{def:g12:deriv:tangent}{tangente}
en $a$. Alors $g$ est \omterm{def:g12:deriv:derivative}{dérivable} et $g'(x) = f'(x) - f'(a)$. Si $f'$ est
\omterm{def:g12:seq:monotonic}{croissante}, $g' \leq 0$ sur $I \cap \intoc{-\infty}{a}$ et $g' \geq 0$ sur
$I \cap \intco{a}{+\infty}$~: $g$ décroît avant $a$ et croît après, donc
$g$ atteint son \omterm{def:g10:functions:extrema}{minimum} $g(a) = 0$, d'où $g \geq 0$.
Les autres implications sont admises à ce niveau~; l'équivalence avec
$f'' \geq 0$ suit du \cref{thm:g12:deriv:variations} appliqué à $f'$.
\end{proof}

\begin{definition}[Point d'inflexion]\label{def:g12:deriv:inflection}
Un point où la courbe de $f$ traverse sa \omterm{def:g12:deriv:tangent}{tangente} est un \emph{point
d'inflexion}\index{point d'inflexion}. Pour $f$ deux fois \omterm{def:g12:deriv:derivative}{dérivable}, les
points d'inflexion sont les points où $f''$ change de signe.
\end{definition}

\begin{example}\label{ex:g12:deriv:cube}
$f(x) = x^3$ a $f''(x) = 6x$~: $f$ est \omterm{def:g12:deriv:convex}{concave} sur $\intoc{-\infty}{0}$,
\omterm{def:g12:deriv:convex}{convexe} sur $\intco{0}{+\infty}$, avec un \omterm{def:g12:deriv:inflection}{point d'inflexion} à l'origine,
où la courbe traverse sa \omterm{def:g12:deriv:tangent}{tangente} (l'axe des $x$).
\end{example}

\begin{omfigure}
\begin{tikzpicture}
\begin{axis}[omaxis,
  width=7.6cm, height=5.6cm,
  xmin=-1.7, xmax=1.7, ymin=-2.9, ymax=2.9,
  xtick={-1,1}, ytick={-1,1}]
  \addplot[omProp, thick, domain=-1.2:1.2] {0};
  \addplot[omDef, thick, domain=-1.4:1.4] {x^3};
  \fill (axis cs:0,0) circle (1.5pt);
  \node[font=\small, anchor=east] at (axis cs:-0.45,-1.7) {concave};
  \node[font=\small, anchor=west] at (axis cs:0.45,1.7) {convexe};
\end{axis}
\end{tikzpicture}

{\small Le \omterm{def:g12:deriv:inflection}{point d'inflexion} de $x \mapsto x^3$ à l'origine~: la courbe
traverse sa \omterm{def:g12:deriv:tangent}{tangente} (orange), \omterm{def:g12:deriv:convex}{concave} à gauche, \omterm{def:g12:deriv:convex}{convexe} à droite.}
\end{omfigure}

\begin{example}[Inégalités de convexité]\label{ex:g12:deriv:ineq}
L'exponentielle est \omterm{def:g12:deriv:convex}{convexe} sur $\R$ (sa \omterm{def:g12:deriv:derivative}{dérivée} seconde est elle-même,
positive). Sa \omterm{def:g12:deriv:tangent}{tangente} en $0$ est $y = 1 + x$, donc
\[
\eu^x \geq 1 + x \quad \text{pour tout } x \in \R .
\]
De telles \emph{inégalités de \omterm{def:g12:deriv:tangent}{tangente}} sont un produit standard de la
\omterm{def:g12:deriv:convex}{convexité}.
\end{example}

\begin{method}[Utiliser la convexité]\label{met:g12:deriv:convexity}
\begin{itemize}
  \item Pour prouver une inégalité $f(x) \geq ax + b$~: identifier la
        droite comme une \omterm{def:g12:deriv:tangent}{tangente} d'une \omterm{def:g11:func:function}{fonction} \omterm{def:g12:deriv:convex}{convexe} $f$ et invoquer le
        \cref{thm:g12:deriv:convexity}(3).
  \item Pour localiser les \omterm{def:g12:deriv:inflection}{points d'inflexion}~: résoudre $f''(x) = 0$
        \emph{et} vérifier que $f''$ change de signe là.
  \item Dans un \omterm{def:g10:functions:table}{tableau de variations}, la \omterm{def:g12:deriv:convex}{convexité} affine le croquis de
        la courbe (de quel côté elle s'infléchit).
\end{itemize}
\end{method}

\section{Exercices}

\begin{exercise}[$\star$]\label{exo:g12:deriv:1}
Dériver les \omterm{def:g11:func:function}{fonctions} suivantes sur leurs \omterm{def:g11:func:function}{ensembles de définition}~:
\[
f(x) = x^3 - 5x + 2, \quad
g(x) = \frac{x}{x^2+1}, \quad
h(x) = (2x+1)^5, \quad
k(x) = \sqrt{x^2 + 1}.
\]
\end{exercise}

\begin{exercise}[$\star$]\label{exo:g12:deriv:2}
Donner l'\omterm{def:g10:algebra:equation}{équation} de la \omterm{def:g12:deriv:tangent}{tangente} à la courbe de $f(x) = x^2 - 3x + 1$ au
point d'\omterm{def:g10:coordgeom:system}{abscisse} $a = 2$, et déterminer les positions de la courbe par
rapport à cette \omterm{def:g12:deriv:tangent}{tangente}.
\end{exercise}

\begin{exercise}[$\star$]\label{exo:g12:deriv:3}
En utilisant la définition de la \omterm{def:g12:deriv:derivative}{dérivée}, calculer
\[
\lim_{x \to 0} \frac{(1+x)^{100} - 1}{x}
\qquad\text{et}\qquad
\lim_{h \to 0} \frac{\sqrt{4+h} - 2}{h}.
\]
\end{exercise}

\begin{exercise}[$\star\star$]\label{exo:g12:deriv:4}
Étudier la \omterm{def:g11:func:function}{fonction} $f(x) = \dfrac{x^2 + 3}{x - 1}$ sur son \omterm{def:g11:func:function}{ensemble de
définition}~: limites, \omterm{def:g12:deriv:derivative}{dérivée}, \omterm{def:g10:functions:table}{tableau de variations}, extrema locaux.
\end{exercise}

\begin{exercise}[$\star\star$]\label{exo:g12:deriv:5}
Une boîte sans couvercle est fabriquée à partir d'une feuille de carton
carrée de côté $30$~cm en découpant des carrés égaux de côté $x$ aux coins
et en relevant les côtés. Déterminer $x$ maximisant le volume de la boîte.
\end{exercise}

\begin{exercise}[$\star\star$]\label{exo:g12:deriv:6}
Soit $f(x) = x^4 - 4x^3 + 10$.
\begin{enumerate}
  \item Calculer $f''$ et déterminer les \omterm{def:g10:numbers:interval}{intervalles} de \omterm{def:g12:deriv:convex}{convexité} et les
        \omterm{def:g12:deriv:inflection}{points d'inflexion} de $f$.
  \item Montrer que la \omterm{def:g12:deriv:tangent}{tangente} au \omterm{def:g12:deriv:inflection}{point d'inflexion} d'\omterm{def:g10:coordgeom:system}{abscisse} $2$ traverse
        la courbe là.
\end{enumerate}
\end{exercise}

\begin{exercise}[$\star\star$]\label{exo:g12:deriv:7}
En utilisant une inégalité de \omterm{def:g12:deriv:tangent}{tangente}, montrer que pour tout $x > 0$,
\[
\sqrt{x} \leq \frac{x + 1}{2},
\]
et identifier quand l'égalité a lieu. (Indication~: la racine carrée est
\omterm{def:g12:deriv:convex}{concave}.)
\end{exercise}

\begin{exercise}[$\star\star\star$]\label{exo:g12:deriv:8}
Soit $f$ \omterm{def:g12:deriv:convex}{convexe} et \omterm{def:g12:deriv:derivative}{dérivable} sur $\R$, et supposons que $f'$ s'annule en
un certain point $a$. Montrer que $f(a)$ est le \omterm{def:g10:functions:extrema}{minimum} \emph{global} de
$f$ sur~$\R$.
\end{exercise}

\begin{exercise}[$\star\star\star$]\label{exo:g12:deriv:9}
Montrer que parmi tous les rectangles de périmètre fixe $p$, le carré a
l'aire la plus grande. Puis montrer que parmi tous les rectangles d'aire
fixe $A$, le carré a le plus petit périmètre, et expliquer comment les
deux énoncés sont liés.
\end{exercise}

\section{Problème~: le calcul du maître nageur}

\begin{problem}[{Devoir du week-end --- le sauvetage le plus rapide
obéit à la loi de Descartes, la convexité certifie tout optimum, et
la canette idéale est exactement aussi haute que large}]
\label{pb:g12:deriv:1}
Un maître nageur aperçoit un nageur en difficulté --- en diagonale
le long de la plage, au large. Courir va plus vite que nager~: la
ligne droite n'est \emph{pas} le trajet le plus rapide. Le trajet qui
minimise le temps se brise au bord de l'eau, et sa \omterm{def:g11:prob:rv}{loi} de brisure
est précisément celle selon laquelle la lumière se réfracte en
entrant dans l'eau~: la nature aussi dérive. Ce problème entraîne la
dérivation des \omterm{def:g11:func:function}{fonctions} composées (\cref{thm:g12:deriv:chain}),
effectue le sauvetage, moissonne les inégalités de la \omterm{def:g12:deriv:convex}{convexité}
(\cref{thm:g12:deriv:convexity}) et dessine une boîte de conserve.

\medskip
\noindent\textbf{Partie I --- Aisance avec les \omterm{def:g11:func:function}{fonctions}
composées.}
\begin{enumerate}
  \item Dériver~: $\left(3x^2 + 1\right)^5$~; \
        $\sqrt{x^2 + 9}$~; \ $\dfrac{1}{x^2 + 1}$.
  \item Donner la \omterm{def:g12:deriv:tangent}{tangente} à $y = \sqrt{x^2 + 9}$ en $x = 4$.
  \item Construire le \omterm{def:g10:functions:table}{tableau de variations} de
        $f(x) = \dfrac{x}{x^2 + 1}$ sur $\R$ (\omterm{def:g12:deriv:extremum}{extremums} compris).
  \item Pour $g(x) = x^3 - 3x^2 + 4$~: \omterm{def:g10:functions:table}{tableau de variations}
        \emph{et} tableau de \omterm{def:g12:deriv:convex}{convexité} ($g''$), avec le \omterm{def:g12:deriv:inflection}{point
        d'inflexion} (\cref{def:g12:deriv:inflection}).
  \item Calculer la \omterm{def:g12:deriv:tangent}{tangente} à $g$ en son \omterm{def:g12:deriv:inflection}{point d'inflexion}, et
        montrer que $g(x) - (\text{tangente}) = (x - 1)^3$~: que dit
        ce changement de signe sur la façon dont la \omterm{def:g12:deriv:tangent}{tangente}
        rencontre la courbe en un \omterm{def:g12:deriv:inflection}{point d'inflexion}~?
\end{enumerate}

\medskip
\noindent\textbf{Partie II --- Le sauvetage.}
Le rivage suit l'axe des \omterm{def:g10:coordgeom:system}{abscisses}~; le maître nageur se tient sur
le sable en $A(0, 30)$, le nageur attend en $B(40, -20)$ (en
mètres). Le sauveteur court à $5$ m/s sur le sable, nage à $2$ m/s,
et entre dans l'eau en un point $(x, 0)$ de son choix.
\begin{enumerate}[resume]
  \item Modélisation~: exprimer le temps total de sauvetage $T(x)$,
        et dire pourquoi seuls les $x \in \intcc{0}{40}$ méritent
        d'être considérés.
  \item Dériver $T$ (la dérivation des composées à l'œuvre --- deux
        fois).
  \item Soit $\theta_1$ l'angle de la portion courue avec la
        \emph{perpendiculaire} au rivage, et $\theta_2$ celui de la
        portion nagée. Montrer que
        $\sin\theta_1 = \frac{x}{\sqrt{x^2 + 900}}$ et que la
        condition $T'(x) = 0$ s'écrit
        \[
        \frac{\sin\theta_1}{5} = \frac{\sin\theta_2}{2}
        \]
        --- la \emph{\omterm{def:g11:prob:rv}{loi} de la réfraction}, avec les deux vitesses à
        la place des vitesses de la lumière dans l'air et dans
        l'eau.
  \item Le segment $[AB]$ coupe le rivage en $x = 24$. Calculer
        $T(24)$, $T(40)$ (courir le long de la plage, puis nager
        droit au large) et $T(0)$~: la ligne droite géométrique
        est-elle la plus rapide~? L'une des deux stratégies
        extrêmes l'est-elle~?
  \item Localiser le point d'entrée optimal par \omterm{thm:g12:limcont:ivt}{dichotomie} sur $T'$
        (il vaut $\approx 33.7$ m)~: donner $x^*$ au mètre près et
        le temps record $T(x^*)$ au dixième de seconde. Combien la
        brisure optimale fait-elle gagner par rapport à la ligne
        droite~?
  \item Pourquoi le point critique est-il un \omterm{def:g10:functions:extrema}{minimum} \emph{global}~?
        (Chaque terme de $T$ est \omterm{def:g12:deriv:convex}{convexe} --- on admettra qu'une
        somme de \omterm{def:g11:func:function}{fonctions} \omterm{def:g12:deriv:convex}{convexes} est \omterm{def:g12:deriv:convex}{convexe} --- puis appliquer
        l'\cref{exo:g12:deriv:8}.)
  \item Le principe de Fermat affirme que la lumière emprunte
        toujours le trajet de \emph{temps} minimal. En déduire
        pourquoi un rayon lumineux se brise exactement à la surface
        en entrant dans l'eau (où la lumière est plus lente), et
        nommer l'observation quotidienne que cela explique à propos
        d'un crayon plongé dans un verre d'eau.
\end{enumerate}

\medskip
\noindent\textbf{Partie III --- La moisson de la \omterm{def:g12:deriv:convex}{convexité}.}
\begin{enumerate}[resume]
  \item Montrer que $x \mapsto x^2$ est \omterm{def:g12:deriv:convex}{convexe}, et démontrer deux
        fois l'inégalité du \omterm{prop:g10:coordgeom:midpoint}{milieu}
        $\left(\frac{a + b}{2}\right)^2 \leq
        \frac{a^2 + b^2}{2}$~: une fois par la \omterm{def:g12:deriv:convex}{convexité}, une fois
        en développant $(a - b)^2 \geq 0$.
  \item En déduire que la \emph{\omterm{def:g11:stat:mean}{moyenne} quadratique}
        $\sqrt{\frac{a^2 + b^2}{2}}$ domine la \omterm{def:g11:stat:mean}{moyenne}
        arithmétique, et assembler la chaîne complète de la série
        --- harmonique $\leq$ géométrique $\leq$ arithmétique
        $\leq$ quadratique --- en vérifiant les quatre pour $a = 2$
        et $b = 8$.
  \item \omterm{def:g12:deriv:tangent}{Tangente} sous la courbe~: $x \mapsto \frac1x$ est \omterm{def:g12:deriv:convex}{convexe}
        sur $\intoo{0}{+\infty}$~; écrire sa \omterm{def:g12:deriv:tangent}{tangente} en $1$ et en
        déduire l'inégalité $\frac1x \geq 2 - x$ pour tout $x > 0$.
        Où y a-t-il égalité~?
  \item Une inflexion dans la vie réelle~: lors d'une épidémie, le
        nombre cumulé de cas suit une courbe en S. Que se passe-t-il
        en son \omterm{def:g12:deriv:inflection}{point d'inflexion}, et pourquoi les épidémiologistes
        guettent-ils cette date~? Relier au signe de la \omterm{def:g12:deriv:derivative}{dérivée}
        seconde de part et d'autre.
\end{enumerate}

\medskip
\noindent\textbf{Partie IV --- La boîte idéale.}
Une boîte cylindrique doit contenir $V = 330$ cm$^3$ avec le moins
de métal possible~: surface $S = 2\pi r^2 + 2\pi r h$, sous la
contrainte $\pi r^2 h = V$.
\begin{enumerate}[resume]
  \item Exprimer $S(r) = 2\pi r^2 + \dfrac{2V}{r}$, dériver, et
        calculer le rayon et la hauteur optimaux pour $V = 330$ (au
        millimètre près).
  \item Démontrer l'élégante \omterm{def:g11:prob:rv}{loi} générale~: à l'optimum, $h = 2r$
        --- la boîte idéale est exactement aussi haute que large.
  \item Vérifier $S''(r) > 0$ et conclure (par
        l'\cref{exo:g12:deriv:8} de nouveau) que l'optimum est
        global. Les canettes de soda réelles sont visiblement plus
        hautes que larges~: donner une raison non mathématique.
  \item Pour finir --- la chaîne de l'optimisation~: modéliser,
        interroger la \omterm{def:g12:deriv:derivative}{dérivée}, résoudre pour trouver les points
        critiques, certifier par la \omterm{def:g12:deriv:convex}{convexité} ou par un tableau,
        interpréter. Dérouler cette liste sur le maître nageur, sur
        la boîte et sur la poutre la plus résistante de l'an dernier
        (\cref{pb:g11:deriv:1}) --- et dire ce que le principe de
        Fermat nous apprend sur qui d'autre déroule cette même
        chaîne.
\end{enumerate}
\end{problem}
